\documentclass[11pt]{article}
\usepackage[margin=0.75in, top=0.4in, bottom=0.5in]{geometry}
\usepackage{listings}
\usepackage{xcolor}
\usepackage{enumitem}
\usepackage{amsmath}
\usepackage{multicol}

\lstset{
  language=Java,
  basicstyle=\ttfamily\small,
  keywordstyle=\bfseries\color{blue!70!black},
  commentstyle=\color{green!50!black},
  stringstyle=\color{red!60!black},
  showstringspaces=false,
  frame=single,
  framesep=4pt,
  xleftmargin=4pt,
  xrightmargin=4pt,
  aboveskip=6pt,
  belowskip=6pt,
}

\pagestyle{empty}

\begin{document}

\begin{center}
  {\Large \textbf{CSCI-UA 102 --- Quiz 6}} \\[4pt]
  {\large Maps and Hash Functions} \\[3pt]
  Name: \underline{\hspace{2.5in}} \quad NetID: \underline{\hspace{1.5in}}
\end{center}

\medskip

\noindent Say I create an implementation of \texttt{PositionList<Integer>} called
\texttt{CrazyList}, and I override \texttt{equals} so that two \texttt{CrazyList}
objects are equal if:
\begin{enumerate}[nosep]
  \item their first elements are equal,
  \item their last elements are equal, and
  \item the sum of all their elements are equal.
\end{enumerate}

\medskip

\noindent\textbf{(1 point)} Write two \texttt{CrazyList}s of length 4 and 5 that are equal.

\bigskip

Length 4: \texttt{[}\underline{\hspace{3.5in}}\texttt{]}

\medskip

Length 5: \texttt{[}\underline{\hspace{3.5in}}\texttt{]}

\bigskip

\noindent\textbf{(4 points)} Write an appropriate implementation of \texttt{hashCode} for \texttt{CrazyList}.

\medskip

\noindent{\footnotesize You may assume these methods:}
\vspace{-0.5em}
\begin{multicols}{2}
\begin{lstlisting}[basicstyle=\ttfamily\scriptsize, aboveskip=3pt, belowskip=3pt]
public interface Position<E> {
    E getElement();
}
\end{lstlisting}
\columnbreak
\begin{lstlisting}[basicstyle=\ttfamily\scriptsize, aboveskip=3pt, belowskip=3pt]
// Methods of PositionList<E>:
Position<E> first();
Position<E> last();
Position<E> after(Position<E> p);
\end{lstlisting}
\end{multicols}

\vfill

\newpage

{\color{blue}
\section*{Solution}

\textbf{Part 1:} Any two lists with the same first element, last element, and sum.

For example: \texttt{[1, 0, 0, 1]} and \texttt{[1, 0, 0, 0, 1]} --- both have first = 1, last = 1, sum = 2.

\medskip

\textbf{Part 2:} Key insight: the \texttt{hashCode} contract requires that
equal objects produce the same hash code. Since \texttt{equals} depends only
on \texttt{first}, \texttt{last}, and \texttt{sum}, the hash code must depend
only on these three values (not on length or individual elements).

\begin{lstlisting}[basicstyle=\ttfamily\small\color{blue},
  keywordstyle=\bfseries\color{blue!70!black},
  commentstyle=\color{blue!50}]
public int hashCode() {
    int sum = 0;
    Position<Integer> current = first();
    while (current != null) {
        sum += current.getElement();
        current = after(current);
    }
    int first = first().getElement();
    int last = last().getElement();
    return 31 * (31 * first + last) + sum;
}
\end{lstlisting}

\textbf{Note:} A simpler \texttt{return first + last + sum;} is also correct ---
it satisfies the contract but produces more collisions.
Any formula using only \texttt{first}, \texttt{last}, and \texttt{sum} is valid.
Using \texttt{length}, individual middle elements, or anything else that
equal lists could disagree on would \textbf{violate} the contract.
}

\end{document}
